\documentclass[11pt,a4paper]{amsart}
\usepackage[margin=1in]{geometry}
\usepackage{amsmath,amssymb,amsthm,mathtools}
\usepackage{booktabs}
\usepackage{enumitem}
\usepackage[colorlinks=true,linkcolor=blue,citecolor=blue,urlcolor=blue]{hyperref}

\newtheorem{theorem}{Theorem}
\newtheorem{lemma}[theorem]{Lemma}
\newtheorem{proposition}[theorem]{Proposition}
\newtheorem{corollary}[theorem]{Corollary}
\theoremstyle{definition}
\newtheorem{definition}[theorem]{Definition}
\newtheorem{remark}[theorem]{Remark}

\newcommand{\Z}{\mathbb Z}
\newcommand{\N}{\mathbb N}

\title[The equation $z^2+y^2z+x^3-2=0$]{An exact divisor parametrisation for the integer solutions of $z^2+y^2z+x^3-2=0$}
\author{Jamal Agbanwa}
\date{\today}

\begin{document}

\begin{abstract}
We give a complete elementary solution of the Diophantine equation
\[
        z^2+y^2z+x^3-2=0,
        \qquad x,y,z\in \Z.
\]
The key point is the factorisation
\[
        z^2+y^2z=z(z+y^2).
\]
Since the case $z=0$ is impossible, every solution is obtained by choosing an integer $x$, a non-zero divisor $z$ of $2-x^3$, and then requiring
\[
        \frac{2-x^3}{z}-z
\]
to be a non-negative square. This gives an exact finite divisor test for each fixed value of $x$, and no external results are required.
\end{abstract}

\maketitle

\section{Statement of the result}

We solve the equation
\begin{equation}\label{eq:main}
        z^2+y^2z+x^3-2=0,
        \qquad x,y,z\in \Z.
\end{equation}
The solution is most naturally expressed as a divisor parametrisation. Throughout the paper, a \emph{non-negative square} means an integer of the form $r^2$ with $r\in\Z$.

\begin{theorem}[Complete divisor parametrisation]\label{thm:main}
The integer solutions of \eqref{eq:main} are exactly the triples
\begin{equation}\label{eq:param_set}
        (x,y,z)\in \Z^3
\end{equation}
with the following three properties:
\begin{enumerate}[label=\textup{(\roman*)}]
    \item $z\ne 0$;
    \item $z\mid (2-x^3)$;
    \item the integer
    \[
            Q(x,z):=\frac{2-x^3}{z}-z
    \]
    is a non-negative square, and $y$ is one of its square roots.
\end{enumerate}
Equivalently,
\begin{equation}\label{eq:solution_set}
\boxed{
\mathcal S
=
\left\{(x,y,z)\in\Z^3:
        z\ne0,
        \ z\mid(2-x^3),
        \ y^2=\frac{2-x^3}{z}-z
\right\}.}
\end{equation}
In particular, for fixed $x$ one obtains the complete finite list of solutions by checking the finitely many divisors $z$ of $2-x^3$.
\end{theorem}

The following equivalent form is often cleaner conceptually.

\begin{theorem}[Factor-pair form]\label{thm:factorpair}
The integer solutions of \eqref{eq:main} are in bijection with the quadruples
\[
        (x,a,b,y)\in\Z^4
\]
satisfying
\begin{equation}\label{eq:factor_conditions}
        ab=2-x^3,
        \qquad b-a=y^2,
        \qquad a\ne0.
\end{equation}
The bijection sends a solution $(x,y,z)$ to
\[
        (x,a,b,y)=(x,z,z+y^2,y),
\]
and the inverse map sends $(x,a,b,y)$ to
\[
        (x,y,z)=(x,y,a).
\]
Equivalently, one may choose a factor pair $(a,b)$ of $2-x^3$ whose difference $b-a$ is a non-negative square, and then set
\[
        z=a,
        \qquad y=\pm\sqrt{b-a},
\]
with the convention that if $b-a=0$, then the single value $y=0$ is obtained.
\end{theorem}

\section{Proof of the parametrisation}

We begin with the only exceptional case that could obstruct division by $z$.

\begin{lemma}\label{lem:znonzero}
If $(x,y,z)\in\Z^3$ satisfies \eqref{eq:main}, then $z\ne0$.
\end{lemma}

\begin{proof}
Suppose, for contradiction, that $z=0$. Substituting into \eqref{eq:main} gives
\[
        x^3-2=0,
\]
so $x^3=2$. But no integer cube is equal to $2$: indeed, if $x\le1$, then $x^3\le1$, while if $x\ge2$, then $x^3\ge8$. This contradiction proves $z\ne0$.
\end{proof}

\begin{lemma}[The fundamental factorisation]\label{lem:factorisation}
For integers $x,y,z$, equation \eqref{eq:main} is equivalent to
\begin{equation}\label{eq:factorised}
        z(z+y^2)=2-x^3.
\end{equation}
\end{lemma}

\begin{proof}
Starting from \eqref{eq:main}, move $x^3-2$ to the other side:
\[
        z^2+y^2z=2-x^3.
\]
The left-hand side factors as
\[
        z^2+y^2z=z(z+y^2).
\]
This gives \eqref{eq:factorised}. Conversely, expanding \eqref{eq:factorised} gives
\[
        z^2+y^2z=2-x^3,
\]
which is exactly \eqref{eq:main}. Hence the two equations are equivalent.
\end{proof}

\begin{lemma}[Divisor-square criterion]\label{lem:criterion}
Let $x,z\in\Z$ with $z\ne0$. There exists an integer $y$ satisfying \eqref{eq:main} if and only if
\[
        z\mid(2-x^3)
\]
and
\[
        \frac{2-x^3}{z}-z
\]
is a non-negative square. In that case the possible values of $y$ are exactly the integer square roots of this number.
\end{lemma}

\begin{proof}
Assume first that there is an integer $y$ such that $(x,y,z)$ satisfies \eqref{eq:main}. By Lemma \ref{lem:factorisation},
\[
        z(z+y^2)=2-x^3.
\]
Because $z\ne0$, this implies immediately that $z$ divides $2-x^3$. Dividing by $z$ gives
\[
        z+y^2=\frac{2-x^3}{z}.
\]
Therefore
\[
        y^2=\frac{2-x^3}{z}-z.
\]
The right-hand side is thus a non-negative square, namely the square of the integer $y$.

Conversely, suppose that $z\mid(2-x^3)$ and that
\[
        \frac{2-x^3}{z}-z=r^2
\]
for some integer $r$. Put $y=r$. Then
\[
        z+y^2
        =z+r^2
        =z+\left(\frac{2-x^3}{z}-z\right)
        =\frac{2-x^3}{z}.
\]
Multiplying both sides by $z$ yields
\[
        z(z+y^2)=2-x^3.
\]
By Lemma \ref{lem:factorisation}, this is equivalent to \eqref{eq:main}. Thus $(x,r,z)$ is a solution. Since the displayed equation forces $y^2=r^2$, the only possible values of $y$ are precisely the integer square roots of $r^2$.
\end{proof}

\begin{proof}[Proof of Theorem \ref{thm:main}]
Let $(x,y,z)$ be an integer solution of \eqref{eq:main}. By Lemma \ref{lem:znonzero}, $z\ne0$. Lemma \ref{lem:criterion} then gives
\[
        z\mid(2-x^3)
        \qquad\text{and}\qquad
        y^2=\frac{2-x^3}{z}-z.
\]
Thus every solution belongs to the set displayed in \eqref{eq:solution_set}.

Conversely, suppose $(x,y,z)$ belongs to the set displayed in \eqref{eq:solution_set}. Then $z\ne0$, $z\mid(2-x^3)$, and
\[
        y^2=\frac{2-x^3}{z}-z.
\]
Rearranging gives
\[
        z+y^2=\frac{2-x^3}{z}.
\]
After multiplication by $z$,
\[
        z(z+y^2)=2-x^3.
\]
By Lemma \ref{lem:factorisation}, this is equivalent to \eqref{eq:main}. Hence $(x,y,z)$ is a solution. This proves both inclusions and therefore proves the complete parametrisation.

Finally, if $x$ is fixed, then $2-x^3\ne0$, because $x^3=2$ has no integer solution. A non-zero integer has only finitely many divisors. Hence the stated procedure checks only finitely many possible values of $z$ and gives the complete finite list for that fixed $x$.
\end{proof}

\begin{proof}[Proof of Theorem \ref{thm:factorpair}]
Let $(x,y,z)$ be a solution of \eqref{eq:main}. Define
\[
        a=z,
        \qquad b=z+y^2.
\]
By Lemma \ref{lem:znonzero}, $a=z\ne0$. By Lemma \ref{lem:factorisation},
\[
        ab=z(z+y^2)=2-x^3.
\]
Moreover,
\[
        b-a=(z+y^2)-z=y^2.
\]
Therefore $(x,a,b,y)$ satisfies \eqref{eq:factor_conditions}.

Conversely, suppose $(x,a,b,y)\in\Z^4$ satisfies \eqref{eq:factor_conditions}. Put $z=a$. Since $b-a=y^2$, we have
\[
        b=z+y^2.
\]
The identity $ab=2-x^3$ becomes
\[
        z(z+y^2)=2-x^3.
\]
By Lemma \ref{lem:factorisation}, $(x,y,z)$ satisfies \eqref{eq:main}.

It remains only to check that the two constructions are inverse to each other. Starting with a solution $(x,y,z)$, the first construction gives
\[
        (x,a,b,y)=(x,z,z+y^2,y),
\]
and the inverse sends this back to $(x,y,z)$. Starting with a quadruple $(x,a,b,y)$ satisfying \eqref{eq:factor_conditions}, the inverse gives $(x,y,a)$; applying the first construction gives
\[
        (x,a,a+y^2,y).
\]
But $b-a=y^2$, so $a+y^2=b$. Hence the original quadruple is recovered. Thus the correspondence is a bijection.
\end{proof}

\section{Consequences and examples}

The equation is even in $y$, and the parametrisation makes this symmetry explicit.

\begin{corollary}[Sign symmetry in $y$]\label{cor:y_symmetry}
If $(x,y,z)$ is an integer solution of \eqref{eq:main}, then $(x,-y,z)$ is also an integer solution. If $y\ne0$, these are distinct solutions; if $y=0$, they coincide.
\end{corollary}

\begin{proof}
The equation contains $y$ only through $y^2$. Therefore replacing $y$ by $-y$ does not change the left-hand side of \eqref{eq:main}. The final assertion is immediate.
\end{proof}

\begin{corollary}[A sign restriction for $x\ge2$]\label{cor:sign_restriction}
If $x\ge2$ and $(x,y,z)$ is a solution of \eqref{eq:main}, then
\[
        z<0<z+y^2.
\]
In particular, $y\ne0$ for every solution with $x\ge2$.
\end{corollary}

\begin{proof}
For $x\ge2$, we have
\[
        2-x^3<0.
\]
By Lemma \ref{lem:factorisation},
\[
        z(z+y^2)=2-x^3<0.
\]
Thus the two factors $z$ and $z+y^2$ have opposite signs. Since $y^2\ge0$, we have
\[
        z\le z+y^2.
\]
The only way two ordered real numbers $z\le z+y^2$ can have opposite signs is
\[
        z<0<z+y^2.
\]
This also implies $y^2>0$, so $y\ne0$.
\end{proof}

For illustration, the following table was obtained directly from Theorem \ref{thm:main}. For each displayed value of $x$, the listed values of $z$ are exactly those shown in the table; each row gives $Q(x,z)=y^2$ and hence the corresponding values of $y$.

\begin{center}
\begin{tabular}{rrrr}
\toprule
$x$ & $2-x^3$ & admissible $z$ & corresponding $y$ \\
\midrule
$1$ & $1$ & $1,-1$ & $0$ \\
$0$ & $2$ & $1,-2$ & $\pm1$ \\
$-2$ & $10$ & $1,-10$ & $\pm3$ \\
$-7$ & $345$ & $5,-69$ & $\pm8$ \\
$8$ & $-510$ & $-15,-34$ & $\pm7$ \\
$20$ & $-7998$ & $-31,-258$ & $\pm17$ \\
$26$ & $-17574$ & $-58,-303$ & $\pm19$ \\
$26$ & $-17574$ & $-87,-202$ & $\pm17$ \\
\bottomrule
\end{tabular}
\end{center}

For example, take $x=8$. Then
\[
        2-x^3=2-512=-510.
\]
The divisor $z=-15$ gives
\[
        Q(8,-15)=\frac{-510}{-15}-(-15)=34+15=49=7^2,
\]
so
\[
        (8,7,-15)
        \quad\text{and}\quad
        (8,-7,-15)
\]
are solutions. Similarly, the divisor $z=-34$ gives
\[
        Q(8,-34)=\frac{-510}{-34}-(-34)=15+34=49=7^2,
\]
so
\[
        (8,7,-34)
        \quad\text{and}\quad
        (8,-7,-34)
\]
are also solutions.

\section{Verification certificate for the sample table}

The table can be checked without any search. Each row follows from the identity
\[
        z^2+y^2z+x^3-2=0.
\]
For instance, the row $(x,y,z)=(26,19,-58)$ gives
\[
        z^2+y^2z+x^3-2
        =(-58)^2+19^2(-58)+26^3-2.
\]
The terms are
\[
        (-58)^2=3364,
        \qquad 19^2(-58)=361(-58)=-20938,
        \qquad 26^3-2=17574.
\]
Therefore
\[
        3364-20938+17574=0.
\]
The other listed rows are verified in the same way, or more systematically by checking the divisor-square identity
\[
        y^2=\frac{2-x^3}{z}-z.
\]
This identity is equivalent to the original equation by Lemma \ref{lem:factorisation}, so the table is not a numerical guess but a direct certificate.

\section{Conclusion}

The factorisation
\[
        z^2+y^2z=z(z+y^2)
\]
reduces the original Diophantine equation exactly to a divisor-square condition. There is no missing exceptional case: $z=0$ would force $x^3=2$, impossible over the integers. Hence every integer solution is obtained by choosing $x\in\Z$, a non-zero divisor $z$ of $2-x^3$, and a square root of
\[
        \frac{2-x^3}{z}-z,
\]
provided this number is a non-negative square. This proves the complete classification of the integer solutions of \eqref{eq:main}.

\end{document}
